\documentclass[10pt,a4paper]{amsart}
\usepackage[margin=19mm]{geometry}
\usepackage{amsmath,amssymb,mathtools}
\usepackage[scr=rsfs,cal=euler]{mathalfa}
\usepackage{xfrac}
\usepackage[unicode,hidelinks,bookmarks=false]{hyperref}
\newcommand{\ie}{i.e.\ }
\newcommand{\cf}{cf.\ }
\newcommand{\resp}{respectively\ }
\newcommand{\Subsection}{\subsection}
\providecommand{\nobreakdash}{\nobreak\hbox{-}}
\setlength{\emergencystretch}{2em}
\multlinegap0pt
\usepackage{amssymb}
\usepackage{mathtools}
\usepackage{bm}
\usepackage[shortlabels]{enumitem}
\usepackage{tikz}
\newtheorem{theorem}{Theorem}
\def \Ga {{\Gamma }}
\title{Cyclic Subgroup Graph of a Group}
\author{Peter J. Cameron, A. Satyanarayana Reddy, Khyati Sharma}
\date{Source: arXiv:2409.13796v2 (CC BY 4.0)}
\begin{document}
\maketitle

\noindent\emph{Source and licence.} Cyclic Subgroup Graph of a Group. Version arXiv:2409.13796v2 (CC BY 4.0), distributed by arXiv under the Creative Commons Attribution 4.0 International licence, http://creativecommons.org/licenses/by/4.0/, as recorded in the arXiv licence metadata for this exact version and shown on the abstract page https://arxiv.org/abs/2409.13796v2. Full licence notice: \texttt{licenses/arxiv-2409.13796-CC-BY-4.0.txt}.\par\smallskip

\noindent\textit{Editorial scope.} Counted unit: the article's theorem t:diambound (source0.tex lines 419--424). The article's complete original proof of it is delivered with it (source0.tex lines 426--433, one \texttt{proof} environment of 8 lines). No mathematics is written, repaired or re-derived: the statement and the proof are the article's own lines, in the article's own notation, and no retained line is substituted or re-flowed.  No further source-local context is needed: every cross-reference of every retained line resolves inside the two delivered spans.  Nothing was omitted from any retained span, so the package carries no omission record.  No retained line contains a citation, so the wrapper carries no bibliography and no bibliography entry.  No retained line contains a figure environment, an image-inclusion command or drawing code, so no image is delivered and the package declares no asset dependency.  Editorial formatting only, in addition: an AMS \texttt{amsart} preamble that restates the article's own declarations and macros used by the retained lines, namely the environments \texttt{theorem} on separate counters, together with the standard packages those lines need.  The retained environments print in this document as Theorem 1 in order of appearance, whereas the article prints the same items with the numbering of its own full document.  The complete native source of the article as distributed in the arXiv e-print for this exact version is bundled unchanged as \texttt{sources/165579/source0.tex}.
\begin{theorem}
Let $G$ be a finite group, and let $g$ be an element of $G$ which maximizes the
number of prime divisors of its order (counted with multiplicity); let this
number be $r$. Then $r\le\Ga(G)\le 2r$.
\label{t:diambound}
\end{theorem}

\begin{proof}
Since one step in the cyclic subgroup graph can only change the number of prime
divisors by $1$, it takes at least $r$ steps to reach $\langle g\rangle$ from
the trivial group; but as we have seen, it can be done in $r$ steps. So the
distance from any cyclic subgroup of $G$ to the trivial subgroup is at most $r$,
and this value is achieved. Moreover, given any two cyclic subgroups $H_1$ and
$H_2$, we can move from either one to the trivial subgroup in at most $r$ steps,so we can move from $H_1$ to $H_2$ in at most $2r$ steps. 
\end{proof}

\end{document}
